% from ega0-20-fr.tex
\oldpage[0\textsubscript{IV-1}]{137}
d'homomorphismes d'anneaux commutatifs. La commutativité du diagramme
correspondant (20.5.7.3) entraîne l'existence d'un unique
$C$-homomorphisme
\begin{equation}
\tag{20.6.4.2}
\label{0.20.6.4.2-fr}
  \Upsilon_{C/B/A}\longrightarrow\Upsilon_{C'/B'/A'}
\end{equation}
déduit canoniquement de (20.6.4.1) et rendant commutatif le diagramme
\begin{equation}
\tag{20.6.4.3}
\label{0.20.6.4.3-fr}
\xymatrix{
  0 \ar[r] & \Upsilon_{C/B/A} \ar[r] \ar[d] &
  \Omega^1_{B/A}\otimes_B C \ar[r]^{v_{C/B/A}} \ar[d] &
  \Omega^1_{C/A} \ar[r]^{u_{C/B/A}} \ar[d] &
  \Omega^1_{C/B} \ar[r] \ar[d] & 0 \\
  0 \ar[r] & \Upsilon_{C'/B'/A'} \ar[r] &
  \Omega^1_{B'/A'}\otimes_{B'} C' \ar[r]_{v_{C'/B'/A'}} &
  \Omega^1_{C'/A'} \ar[r]_{u_{C'/B'/A'}} &
  \Omega^1_{C'/B'} \ar[r] & 0
}
\end{equation}

La donnée de l'homomorphisme (20.6.4.2) équivaut d'ailleurs à celle d'un
$C'$-homomorphisme
\begin{equation}
\tag{20.6.4.4}
\label{0.20.6.4.4-fr}
  \Upsilon_{C/B/A}\otimes_C C'\longrightarrow\Upsilon_{C'/B'/A'}
\end{equation}
qui, composé avec l'homomorphisme canonique
$\Upsilon_{C/B/A}\to\Upsilon_{C/B/A}\otimes_C C'$, redonne (20.6.4.2).
Il est clair que (20.6.4.2) jouit d'une propriété évidente de
transitivité, permettant de dire que $\Upsilon_{C/B/A}$ est un
\emph{foncteur} en le triplet $(A,B,C)$.
\end{env}

\begin{env}[20.6.5]
\label{0.20.6.5-fr}
Il nous sera commode pour la suite, sous les conditions de (20.6.1),
d'introduire un \emph{complexe} (de chaînes) \emph{de $C$-modules}
$\mathrm{K}_\bullet(C/B/A)$ dont les termes sont nuls sauf en degrés 0
et 1, où on prend
\begin{equation}
\tag{20.6.5.1}
\label{0.20.6.5.1-fr}
\left\{
\begin{aligned}
  \mathrm{K}_0(C/B/A)&=\Omega^1_{C/A},\\
  \mathrm{K}_1(C/B/A)&=\Omega^1_{B/A}\otimes_B C,
\end{aligned}
\right.
\end{equation}
l'opérateur différentiel $\mathrm{K}_1\to\mathrm{K}_0$ étant
$v_{C/B/A}$. Cela permet d'écrire (à des isomorphismes canoniques près)
$\Omega^1_{C/B}$ et $\Upsilon_{C/B/A}$ comme les \emph{modules
d'homologie} de ce complexe
\begin{equation}
\tag{20.6.5.2}
\label{0.20.6.5.2-fr}
  \mathrm{H}_0(\mathrm{K}_\bullet(C/B/A))=\Omega^1_{C/B},\qquad
  \mathrm{H}_1(\mathrm{K}_\bullet(C/B/A))=\Upsilon_{C/B/A}.
\end{equation}
\end{env}

De même :
\begin{proposition}[20.6.6]
\label{0.20.6.6-fr}
Sous les hypothèses de (20.6.1), pour tout $C$-module $L$, on a des
$C$-isomorphismes canoniques
\begin{align}
\tag{20.6.6.1}
\label{0.20.6.6.1-fr}
  \mathrm{H}^0(\mathrm{K}_\bullet(C/B/A),L)&\simeq\operatorname{D\acute{e}r}_B(C,L),\\
\tag{20.6.6.2}
\label{0.20.6.6.2-fr}
  \mathrm{H}^1(\mathrm{K}_\bullet(C/B/A),L)&\simeq
  \operatorname{Exalcom}_{B/A}(C,L).
\end{align}
\end{proposition}

\begin{proof}
En effet, le complexe de cochaînes
$\operatorname{Hom}_C(\mathrm{K}_\bullet(C/B/A),L)$ n'est autre, en vertu de (20.4.8)
et (20.5.6), que le complexe
\[
  \cdots\longrightarrow0\longrightarrow\operatorname{D\acute{e}r}_A(C,L)
  \longrightarrow\operatorname{D\acute{e}r}_A(B,L)\longrightarrow0\longrightarrow\cdots
\]
